\section{Background}
\label{sec:background}

% This paper studies context parameterization under continual updates, where contextual information is internalized into a parameterized state for reuse across queries. 
% Let $C$, $q$, and $y$ denote the context, query, and answer sequences, respectively, and let $p_\theta$ denote the conditional generation distribution of a language model with parameters $\theta$.
% Conventional long-context and RAG methods instead supply context explicitly at inference time, generating answers via $p_\theta(y\mid C,q)$~\citep{lewis2020retrieval,guu2020retrieval,ding2023longnet}.
% This incurs repeated reprocessing overhead per query, and long-context models may further suffer degraded information access depending on context position~\citep{liu2024lost,hsieh2024found}.
We study context parameterization under continual updates, internalizing information into a parameterized state reused across queries as $C$ grows monotonically.
Let $C$, $q$, and $y$ denote the context, query, and answer, and $p_\theta(\cdot\mid\cdot)$ the conditional distribution of a $\theta$-parameterized language model.
Conventional long-context and RAG methods supply context at inference time---concatenating context into the input window or retrieving passages---via $p_\theta(y\mid C,q)$~\citep{lewis2020retrieval,guu2020retrieval,ding2023longnet}, incurring reprocessing costs scaling with $|C|$, compounding across queries~\citep{gao2023retrieval}, and suffering position bias, underutilizing evidence away from the input's boundaries~\citep{liu2024lost,hsieh2024found}.
These costs motivate context parameterization internalizing $C$ into parameters decoupled from per-query inference.

Context Distillation (CD) instead transfers contextual information from explicit input into model parameters.
Given a context $C$ and an associated query distribution $\mathcal{Q}_C$, CD uses a context-conditioned teacher to optimize a context-specific parameter state $\theta_C$:
\begin{equation}
\min_{\theta_C}\mathbb{E}_{q\sim\mathcal{Q}_C}\left[
D_{\mathrm{KL}}\left(p_\theta(\cdot\mid C,q)\parallel
p_{\theta_C}(\cdot\mid q)\right)\right].
\end{equation}
The teacher has direct access to $C$, whereas the student is optimized to reproduce the teacher distribution using only the query.
In this way, information originally supplied through the input context is encoded into the resulting parameter state.
After distillation, the model can answer queries without explicitly receiving $C$ during inference, allowing the same parameterized context to be reused across multiple queries~\citep{askell2021general,snell2022learning}.
However, CD requires context-specific distillation data and iterative optimization for each new context, making repeated updates costly.

Doc-to-LoRA (D2L) amortizes this per-context optimization using a hypernetwork $H_\phi$ that maps a context directly to a LoRA weight delta~\citep{hu2022lora,charakorn2026doc}:
\begin{equation}
\Delta W_C=H_\phi(C),\qquad
\theta_C=\theta\oplus\Delta W_C .
\end{equation}
Here, $\Delta W_C$ denotes the generated LoRA adapter and $\oplus$ denotes applying it to the frozen base model.
Instead of optimizing a separate parameter state for each context, D2L learns the mapping from contextual information to parameter updates across a collection of training contexts.
Once trained, $H_\phi$ generates a reusable context-specific adapter in a single forward pass, substantially reducing the parameterization cost for previously unseen contexts.
The generated adapter can be independently applied to or removed from the frozen base model, providing a lightweight parameter representation of the corresponding context.
However, D2L parameterizes contexts without explicitly modeling continual updates to existing parameterized states.